\section{Introduction}

Orange routes origin--destination demands on a directed IP/MPLS backbone.
Between routing changes, traffic follows shortest paths induced by administrative metrics,
and equal-cost paths share the flow.
Segment Routing~\citep{rfc8402} inserts a short list of waypoints; between successive
waypoints the network still uses shortest-path forwarding, so a compact header can move
traffic away from a congested link.

The EURO/ROADEF 2026 challenge~\citep{orange2026} defines $T$-Adaptive Segment Routing
($T$-ASR) over a general horizon $T=\{0,\ldots,h-1\}$.
In the course experiments we take $h=2$: period~0 is the nominal topology, and period~1
withdraws some links for maintenance while the traffic matrix changes.
A budget limits how many segments may change between the two periods.

The challenge objective~\citep{orange2026} sorts every arc-period load and minimizes that
vector lexicographically.
This report solves only the primary scalar stage of that order: the maximum link utilization.
Section~\ref{sec:obj} states the scientific rationale for that choice.
We solve all twenty Set~A instances on period~0 with a limit of 1800\,s, then on both periods
with the budget and a limit of 3600\,s.
Every table entry is a mean over ten runs with seeds 2026--2035.
